\uuid{Dqaj}
\exo7id{1631}
\titre{exo7 1631}
\auteur{barraud}
\organisation{exo7}
\datecreate{2003-09-01}
\isIndication{false}
\isCorrection{false}
\chapitre{Réduction d'endomorphisme, polynôme annulateur}
\sousChapitre{Diagonalisation}
\cpge{mpsi-mp2i.polynomes-fractions-rationnelles.b, pcsi.polynomes.b, ptsi.polynomes.b, mpsi-mp2i.applications-lineaires.a, pcsi.applications-lineaires.a, ptsi.applications-lineaires.a, mpsi-mp2i.applications-lineaires.b, pcsi.applications-lineaires.b, ptsi.applications-lineaires.b, bcpst1.applications-lineaires-matrices, mp-mpi.reduction-endomorphismes.d, pc.reduction.c, psi.reduction-endomorphismes.c, pt.reduction-endomorphismes.c, mp-mpi.reduction-endomorphismes.i, pc.reduction.d, psi.reduction-endomorphismes.d, bcpst2.valeurs-propres-vecteurs-propres.b}
\module{Algèbre}
\niveau{L2}
\difficulte{}

\contenu{
\question{
Soit $P_{0}$ un polynôme de $\R_{n}[X]$, et $f$ l'application
suivante :
$$
 f: \begin{array}{rcl}
      \R_{n}[X] & \rightarrow  & \R_{n}[X] \\[1ex]
      P         & \mapsto & R=
      \text{ reste de la division euclidienne de $P$ par $P_{0}$}
    \end{array}
$$
A l'aide d'un polynôme annulateur de $f$, montrer que $f$ est
diagonalisable.
}
}
